\section{Anthropic 方向的启示}

\subsection{约束式推理哲学}

Anthropic 围绕 Claude 3.7 和 Claude 4 的公开表述相对克制：强调集成推理、用户控制的 thinking budget、真实世界任务、代码质量，以及后来在扩展思考中使用工具的能力。

\begin{importantbox}{更长的推理链 $\neq$ 更高的智能}
\textit{Producing a longer reasoning trace doesn't automatically make a model more intelligent.}（产生更长的推理轨迹并不会自动让模型更智能。）过度的可见推理往往意味着分配失败——模型未能优先化、压缩或行动。Anthropic 的路径暗示了一种更有纪律的观点：thinking 应该由目标工作负载来塑造。
\end{importantbox}

\subsection{从训练模型到训练 Agent}

这种对目标效用的强调指向了更大的图景。正如 Qwen3 博客中所写：\textit{``We are transitioning from an era focused on training models to one centered on training agents.''}（我们正在从专注于训练模型的时代过渡到以训练智能体为中心的时代。）

\begin{knowledgebox}{Agent 的定义}
Agent 是一个能够制定计划、决定何时行动、使用工具、感知环境反馈、修正策略，并在长时间跨度内持续运行的系统。它由与世界的\textbf{闭环交互}（closed-loop interaction）来定义。
\end{knowledgebox}

\subsection{本章小结}

Anthropic 的克制风格提供了有价值的纠正：推理应服务于具体工作负载而非炫技。更长的推理链不等于更高的智能。整个行业正从``训练模型''向``训练 Agent''转变。


